\section{Classification of Coxeter Graphs}

Let \(\Gamma\) be the Coxeter graph associated with a set of simple roots.
We normalize the corresponding vectors \(v_i\) so that
\[
    (v_i,v_i)=1.
\]
Recall that, from the previous discussion, for every vertex \(v_i\), the
number of edges adjacent to \(v_i\), counted with multiplicity, is at most
\(4\).

We shall classify the possible connected Coxeter graphs. The following
reductions are used.

\begin{enumerate}
    \item If a triple edge occurs, then the graph must be the whole graph:
          \[
              \begin{tikzpicture}[coxeter]
                  \coordinate (a) at (0,0);
                  \coordinate (b) at (1.2,0);
                  \draw (a) -- (b);
                  \draw ([yshift=2pt]a) -- ([yshift=2pt]b);
                  \draw ([yshift=-2pt]a) -- ([yshift=-2pt]b);
                  \node[vertex] at (a) {};
                  \node[vertex] at (b) {};
              \end{tikzpicture}
          \]
          This gives the graph of type \(G_2\).

    \item If a double edge occurs, then the graph is a chain with a double
          edge somewhere in the chain. Thus it has the form
          \[
              \begin{tikzpicture}[coxeter]
                  \coordinate (v1) at (0,0);
                  \coordinate (v2) at (0.8,0);
                  \coordinate (vp) at (2.1,0);
                  \coordinate (wq) at (3.2,0);
                  \coordinate (w2) at (4.5,0);
                  \coordinate (w1) at (5.3,0);

                  \draw (v1) -- (v2);
                  \node at (1.45,0) {\(\cdots\)};
                  \draw (1.75,0) -- (vp);

                  \draw[double, double distance=2pt] (vp) -- (wq);

                  \draw (wq) -- (3.65,0);
                  \node at (4.0,0) {\(\cdots\)};
                  \draw (w2) -- (w1);

                  \node[vertex] at (v1) {};
                  \node[vertex] at (v2) {};
                  \node[vertex] at (vp) {};
                  \node[vertex] at (wq) {};
                  \node[vertex] at (w2) {};
                  \node[vertex] at (w1) {};

                  \node at (v1) [below=4pt] {\(v_1\)};
                  \node at (v2) [below=4pt] {\(v_2\)};
                  \node at (vp) [below=4pt] {\(v_p\)};
                  \node at (wq) [below=4pt] {\(w_q\)};
                  \node at (w2) [below=4pt] {\(w_2\)};
                  \node at (w1) [below=4pt] {\(w_1\)};
              \end{tikzpicture}
          \]

    \item If all edges are single, then the graph is either a chain, or a
          tree with one branching point of degree \(3\).
\end{enumerate}

\subsection{The case of a double edge}

Assume that the graph contains a double edge. Let the two chains attached
to the double edge have lengths \(p\) and \(q\), respectively. Define
\[
    v=v_1+2v_2+\cdots+p v_p,
    \qquad
    w=w_1+2w_2+\cdots+q w_q.
\]
Since \(v\) and \(w\) are linearly independent, by the strict Cauchy--Schwarz
inequality,
\[
    (v,w)^2 < (v,v)(w,w).
\]

We compute:
\[
    \begin{aligned}
        (v,v)
         & =
        1^2+2^2+\cdots+p^2
        -
        2\cdot \frac12
        \bigl(1\cdot 2+2\cdot 3+\cdots+(p-1)p\bigr) \\
         & =
        p^2-\frac{p(p-1)}{2}
        =
        \frac{p(p+1)}{2}.
    \end{aligned}
\]
Similarly,
\[
    (w,w)=\frac{q(q+1)}{2}.
\]
Since \(v_p\) and \(w_q\) are joined by a double edge,
\[
    (v_p,w_q)=-\frac{\sqrt{2}}{2}.
\]
Hence
\[
    (v,w)=-pq\frac{\sqrt{2}}{2},
\]
and therefore
\[
    (v,w)^2=\frac{p^2q^2}{2}.
\]

Thus
\[
    \frac{p^2q^2}{2}
    <
    \frac{p(p+1)}{2}\cdot \frac{q(q+1)}{2}.
\]
Cancelling \(pq>0\), we obtain
\[
    2pq < (p+1)(q+1).
\]
Equivalently,
\[
    pq-p-q+1<2,
\]
that is,
\[
    (p-1)(q-1)<2.
\]
Hence either \(p=1\), or \(q=1\), or
\[
    p=q=2.
\]
The first two possibilities give the Coxeter graph of type \(B_n\) or
\(C_n\), while the last possibility gives the Coxeter graph of type \(F_4\).

\subsection{The simply-laced branching case}

Now assume that all edges are single and that the graph has a branching
point. Let the branching vertex be \(\gamma\), and let the three arms have
lengths \(p,q,r\). Write
\[
    v=v_1+2v_2+\cdots+p v_p,
\]
\[
    w=w_1+2w_2+\cdots+q w_q,
\]
\[
    u=u_1+2u_2+\cdots+r u_r.
\]

The vectors \(v,w,u\) are pairwise orthogonal. Extend
\[
    \frac{v}{\sqrt{(v,v)}},
    \qquad
    \frac{w}{\sqrt{(w,w)}},
    \qquad
    \frac{u}{\sqrt{(u,u)}}
\]
to an orthonormal basis. Since \((\gamma,\gamma)=1\), we have
\[
    \frac{(\gamma,v)^2}{(v,v)}
    +
    \frac{(\gamma,w)^2}{(w,w)}
    +
    \frac{(\gamma,u)^2}{(u,u)}
    < 1.
\]
Here the inequality is strict because \(\gamma,v,w,u\) are linearly
independent.

As before,
\[
    (v,v)=\frac{p(p+1)}{2},
    \qquad
    (w,w)=\frac{q(q+1)}{2},
    \qquad
    (u,u)=\frac{r(r+1)}{2}.
\]
Since \(\gamma\) is joined to \(v_p,w_q,u_r\) by single edges,
\[
    (\gamma,v)=-\frac{p}{2},
    \qquad
    (\gamma,w)=-\frac{q}{2},
    \qquad
    (\gamma,u)=-\frac{r}{2}.
\]
Therefore
\[
    \frac{p}{2(p+1)}
    +
    \frac{q}{2(q+1)}
    +
    \frac{r}{2(r+1)}
    <1.
\]
Equivalently,
\[
    \frac{1}{p+1}
    +
    \frac{1}{q+1}
    +
    \frac{1}{r+1}
    >1.
\]

Assume, without loss of generality, that
\[
    p\le q\le r.
\]
Then
\[
    \frac{3}{p+1}>1,
\]
so \(p=1\).

Now
\[
    \frac12+\frac{1}{q+1}+\frac{1}{r+1}>1.
\]
Hence
\[
    \frac{1}{q+1}+\frac{1}{r+1}>\frac12.
\]
Since \(q\le r\), this forces \(q=1\) or \(q=2\).

\begin{itemize}
    \item If \(q=1\), then \(r\) may be arbitrary. This gives the family
          \(D_n\).

    \item If \(q=2\), then
          \[
              \frac12+\frac13+\frac{1}{r+1}>1,
          \]
          so
          \[
              \frac{1}{r+1}>\frac16.
          \]
          Hence
          \[
              r+1<6,
          \]
          and therefore
          \[
              r=2,3,4.
          \]
          These give the exceptional types
          \[
              E_6,\qquad E_7,\qquad E_8.
          \]
\end{itemize}

\subsection{Complete classification}

Combining the preceding cases, the connected Coxeter graphs are precisely
the following.

\[
    \begin{array}{ccl}
        G_2
         &
        \begin{tikzpicture}[coxeter]
            \coordinate (a) at (0,0);
            \coordinate (b) at (1.2,0);
            \draw (a) -- (b);
            \draw ([yshift=2pt]a) -- ([yshift=2pt]b);
            \draw ([yshift=-2pt]a) -- ([yshift=-2pt]b);
            \node[vertex] at (a) {};
            \node[vertex] at (b) {};
        \end{tikzpicture}
         &
        \\[1.5em]

        A_n
         &
        \begin{tikzpicture}[coxeter]
            \foreach \x in {0,1,2,4,5}
            \node[vertex] at (\x,0) {};
            \draw (0,0)--(1,0)--(2,0);
            \node at (3,0) {\(\cdots\)};
            \draw (4,0)--(5,0);
        \end{tikzpicture}
         &
        \\[1.5em]

        B_n \text{ or } C_n
         &
        \begin{tikzpicture}[coxeter]
            \foreach \x in {0,1,2,4,5}
            \node[vertex] at (\x,0) {};
            \draw (0,0)--(1,0)--(2,0);
            \node at (3,0) {\(\cdots\)};
            \draw[double, double distance=2pt] (4,0)--(5,0);
        \end{tikzpicture}
         &
        \\[1.5em]

        F_4
         &
        \begin{tikzpicture}[coxeter]
            \foreach \x in {0,1,2,3}
            \node[vertex] at (\x,0) {};
            \draw (0,0)--(1,0);
            \draw[double, double distance=2pt] (1,0)--(2,0);
            \draw (2,0)--(3,0);
        \end{tikzpicture}
         &
        \\[1.5em]

        D_n
         &
        \begin{tikzpicture}[coxeter]
            \node[vertex] (c) at (1,0) {};
            \node[vertex] (a) at (0,0.7) {};
            \node[vertex] (b) at (0,-0.7) {};
            \node[vertex] (d) at (2,0) {};
            \node[vertex] (e) at (3,0) {};
            \node[vertex] (f) at (5,0) {};
            \draw (c)--(a);
            \draw (c)--(b);
            \draw (c)--(d)--(e);
            \node at (4,0) {\(\cdots\)};
            \draw (4.35,0)--(f);
        \end{tikzpicture}
         &
        \\[2em]

        E_6
         &
        \begin{tikzpicture}[coxeter]
            \foreach \x in {0,1,2,3,4}
            \node[vertex] at (\x,0) {};
            \node[vertex] at (2,-0.8) {};
            \draw (0,0)--(1,0)--(2,0)--(3,0)--(4,0);
            \draw (2,0)--(2,-0.8);
        \end{tikzpicture}
         &
        \\[2em]

        E_7
         &
        \begin{tikzpicture}[coxeter]
            \foreach \x in {0,1,2,3,4,5}
            \node[vertex] at (\x,0) {};
            \node[vertex] at (2,-0.8) {};
            \draw (0,0)--(1,0)--(2,0)--(3,0)--(4,0)--(5,0);
            \draw (2,0)--(2,-0.8);
        \end{tikzpicture}
         &
        \\[2em]

        E_8
         &
        \begin{tikzpicture}[coxeter]
            \foreach \x in {0,1,2,3,4,5,6}
            \node[vertex] at (\x,0) {};
            \node[vertex] at (2,-0.8) {};
            \draw (0,0)--(1,0)--(2,0)--(3,0)--(4,0)--(5,0)--(6,0);
            \draw (2,0)--(2,-0.8);
        \end{tikzpicture}
         &
    \end{array}
\]

\begin{definition}[Dynkin diagram]
    Let \(\Phi\) be a root system. Whenever a double or triple edge occurs in
    the Coxeter graph of \(\Phi\), we add an arrow pointing to the shorter of the
    two roots. The resulting oriented graph is called the \emph{Dynkin diagram}
    of \(\Phi\).
\end{definition}

Thus the Coxeter graphs of \(B_n\) and \(C_n\) have the same underlying
unoriented graph, but their Dynkin diagrams are distinguished by the
direction of the arrow:
\[
    B_n:
    \qquad
    \begin{tikzpicture}[coxeter]
        \foreach \x in {0,1,2,4,5}
        \node[vertex] at (\x,0) {};
        \draw (0,0)--(1,0)--(2,0);
        \node at (3,0) {\(\cdots\)};
        \draw[double, double distance=2pt, dynarrow] (4,0)--(5,0);
    \end{tikzpicture}
\]
\[
    C_n:
    \qquad
    \begin{tikzpicture}[coxeter]
        \foreach \x in {0,1,2,4,5}
        \node[vertex] at (\x,0) {};
        \draw (0,0)--(1,0)--(2,0);
        \node at (3,0) {\(\cdots\)};
        \draw[double, double distance=2pt, dynarrow] (5,0)--(4,0);
    \end{tikzpicture}
\]

\tikzset{
rootdiagram/.style={baseline=-0.5ex, scale=0.9},
vertex/.style={circle, fill=black, inner sep=1.4pt},
dynarrow/.style={-{Latex[length=2mm,width=1.5mm]}},
dynarrowleft/.style={{Latex[length=2mm,width=1.5mm]}-}
}

